
\begin{propbox}
\bb{Règle des signes} :
    \begin{itemize}
        \item Le produit de deux nombres de \bb{même signe} est \answer{positif} ;
        \item Le produit de deux nombres de \bb{signes différents} est
            \answer{négatif}.
    \end{itemize}
    \begin{center}
        \begin{tabular}{c|c|c}
        Signe & Signe & Résultat \\ \hline
        $+$ & $+$ & $+$ \\
        $-$ & $-$ & $+$ \\
        $+$ & $-$ & $-$ \\
        $-$ & $+$ & $-$
        \end{tabular}
    \end{center}
\end{propbox}

\begin{propbox}
    Pour \bb{multiplier deux nombres relatifs}, 
    \begin{enumerate}
        \item On multiplie leurs leur distance à zéro (valeurs absolues).
        \item On applique la règle des signes. 
    \end{enumerate}
\end{propbox}


\begin{exemplebox}
    \begin{align*}
        (+4)\times(+3)&= (\answer{+\times+})\times4\times3 = \answer{+12} \\ 
        (+4)\times(-3)&= (\answer{+\times-})\times4\times3 = \answer{-12} \\
        (-4)\times(-3)&= (\answer{-\times-})\times4\times3 = \answer{+12} \\
        (-4)\times(+3)&= (\answer{-\times+})\times4\times3 = \answer{-12} 
    \end{align*}
\end{exemplebox}


\begin{attentionbox}
\begin{itemize}
    \item Deux signes \bb{identiques} donnent un résultat \answer{positif}.
    \item Deux signes \bb{différents} donnent un résultat \answer{négatif}.
\end{itemize}
\end{attentionbox}

\begin{propbox}
    Le produit de plusieurs nombres relatifs est 
    \begin{itemize}
        \item \bb{positif} s'il comprend un nombre \bb{pair} de facteurs
            négatifs.
        \item \bb{négatif} s'il comprend un nombre \bb{impair} de facteurs négatifs.
    \end{itemize}
\end{propbox}
